% GENERATED FILE. Edit canonical lessons, then run npm run generate:books.
% canonical-source-hash: fnv1a64:0abb4431fdabddd6
% canonical-lessons: SA-C234-susyati, SA-C234-plavate, SA-C234-arabhate, SA-C234-samapayati, SA-C234-prayacchati

\chapter{Dries up, Floats, Begins, Finishes, Hands over}
\label{ch:sa-dries-up-floats-begins-finishes-hands-over}
\hlchaptermodality{234}

\begin{chapteropening}
\textbf{By the end of this chapter:} \emph{I can say dries up, floats, begins, finishes and hands over.}

Complete the last lesson of chapter 234: I can say dries up, floats, begins, finishes and hands over.
\end{chapteropening}

\section[\texorpdfstring{\sk{शुष्यति} (śuṣyati)}{शुष्यति (śuṣyati)}]{\sk{शुष्यति} (śuṣyati) — dries up}

\label{lesson:SA-C234-susyati}



\begin{quote}
Before the new one: say the Sanskrit for follows, then the Sanskrit for goes out.
\end{quote}

\subsection*{You'll want to know: \sk{शुष्यति}}
\textbf{\sk{शुष्यति}} — \emph{śuṣyati} — \textquotedblleft{}dries up\textquotedblright{}.

\begin{cousinweb}
\emph{śuṣyati} is the present of the root √śuṣ, \textquotedblleft{}to dry, wither\textquotedblright{} (MW). MW thinks the root was first \emph{suṣ}, and compares Greek \emph{auō}, \textquotedblleft{}I dry\textquotedblright{}, and Lithuanian \emph{sausas}, \textquotedblleft{}dry\textquotedblright{}, so it is inherited from Indo-European.
\end{cousinweb}

\begin{grammarlens}[title={What you've built}]
One more word for this chapter.
\end{grammarlens}

\subsection*{Your turn}
\begin{itemize}
\raggedright
  \item \emph{Say it:} \emph{śuṣyati}
  \item \emph{Say it:} \emph{śuṣyati}, once more
  \item \emph{Say it:} say \emph{śuṣyati}
  \item \emph{Recall:} say the Sanskrit for follows, then the Sanskrit for goes out, then say \emph{śuṣyati} again
\end{itemize}

\begin{tcolorbox}[breakable,colback=teal!4,colframe=teal!35!black,title={Before you move on}]
What does \sk{शुष्यति} mean? (\textquotedblleft{}Dries up\textquotedblright{}.) Say it once more.
\end{tcolorbox}
\section[\texorpdfstring{\sk{प्लवते} (plavate)}{प्लवते (plavate)}]{\sk{प्लवते} (plavate) — floats}

\label{lesson:SA-C234-plavate}



\begin{quote}
Before the new one: say the Sanskrit for goes out, then the Sanskrit for dries up.
\end{quote}

\subsection*{You'll want to know: \sk{प्लवते}}
\textbf{\sk{प्लवते}} — \emph{plavate} — \textquotedblleft{}floats\textquotedblright{}.

\begin{cousinweb}
\emph{plavate} is the present of the root √plu, \textquotedblleft{}to float, swim\textquotedblright{} (MW), a Rigvedic verb (CDIAL 9024). The root is inherited from Indo-European: MW compares Greek \emph{pleō}, \textquotedblleft{}I sail\textquotedblright{}, and Latin \emph{pluit}, \textquotedblleft{}it rains\textquotedblright{}.
\end{cousinweb}

\begin{grammarlens}[title={What you've built}]
One more word for this chapter.
\end{grammarlens}

\subsection*{Your turn}
\begin{itemize}
\raggedright
  \item \emph{Say it:} \emph{plavate}
  \item \emph{Say it:} \emph{plavate}, once more
  \item \emph{Say it:} say \emph{plavate}
  \item \emph{Recall:} say the Sanskrit for goes out, then the Sanskrit for dries up, then say \emph{plavate} again
\end{itemize}

\begin{tcolorbox}[breakable,colback=teal!4,colframe=teal!35!black,title={Before you move on}]
What does \sk{प्लवते} mean? (\textquotedblleft{}Floats\textquotedblright{}.) Say it once more.
\end{tcolorbox}
\section[\texorpdfstring{\sk{आरभते} (ārabhate)}{आरभते (ārabhate)}]{\sk{आरभते} (ārabhate) — begins}

\label{lesson:SA-C234-arabhate}



\begin{quote}
Before the new one: say the Sanskrit for dries up, then the Sanskrit for floats.
\end{quote}

\subsection*{You'll want to know: \sk{आरभते}}
\textbf{\sk{आरभते}} — \emph{ārabhate} — \textquotedblleft{}begins\textquotedblright{}.

\begin{cousinweb}
\emph{ārabhate} is built on \emph{ā-} and the root √rabh, \textquotedblleft{}to take hold of\textquotedblright{} (MW), which MW connects with √labh, \textquotedblleft{}to obtain\textquotedblright{}. In the Rigveda it means \textquotedblleft{}seizes\textquotedblright{}, and from the Śatapatha Brāhmaṇa \textquotedblleft{}undertakes, begins\textquotedblright{} (CDIAL 1305): to begin is to take hold of a task.
\end{cousinweb}

\begin{grammarlens}[title={What you've built}]
One more word for this chapter.
\end{grammarlens}

\subsection*{Your turn}
\begin{itemize}
\raggedright
  \item \emph{Say it:} \emph{ārabhate}
  \item \emph{Say it:} \emph{ārabhate}, once more
  \item \emph{Say it:} say \emph{ārabhate}
  \item \emph{Recall:} say the Sanskrit for dries up, then the Sanskrit for floats, then say \emph{ārabhate} again
\end{itemize}

\begin{tcolorbox}[breakable,colback=teal!4,colframe=teal!35!black,title={Before you move on}]
What does \sk{आरभते} mean? (\textquotedblleft{}Begins\textquotedblright{}.) Say it once more.
\end{tcolorbox}
\section[\texorpdfstring{\sk{समापयति} (samāpayati)}{समापयति (samāpayati)}]{\sk{समापयति} (samāpayati) — finishes}

\label{lesson:SA-C234-samapayati}



\begin{quote}
Before the new one: say the Sanskrit for floats, then the Sanskrit for begins.
\end{quote}

\subsection*{You'll want to know: \sk{समापयति}}
\textbf{\sk{समापयति}} — \emph{samāpayati} — \textquotedblleft{}finishes\textquotedblright{}.

\begin{cousinweb}
\emph{samāpayati} is the causative of \emph{sam-āp}, \textquotedblleft{}to obtain completely\textquotedblright{}, from the root √āp, \textquotedblleft{}to reach, obtain\textquotedblright{} (MW). It first meant \textquotedblleft{}causes to gain completely\textquotedblright{} and then \textquotedblleft{}finishes\textquotedblright{} (CDIAL 13216). MW compares √āp with Latin \emph{apīscor}, \textquotedblleft{}I reach\textquotedblright{}, and \emph{aptus}, \textquotedblleft{}fitted\textquotedblright{}.
\end{cousinweb}

\begin{grammarlens}[title={What you've built}]
One more word for this chapter.
\end{grammarlens}

\subsection*{Your turn}
\begin{itemize}
\raggedright
  \item \emph{Say it:} \emph{samāpayati}
  \item \emph{Say it:} \emph{samāpayati}, once more
  \item \emph{Say it:} say \emph{samāpayati}
  \item \emph{Recall:} say the Sanskrit for floats, then the Sanskrit for begins, then say \emph{samāpayati} again
\end{itemize}

\begin{tcolorbox}[breakable,colback=teal!4,colframe=teal!35!black,title={Before you move on}]
What does \sk{समापयति} mean? (\textquotedblleft{}Finishes\textquotedblright{}.) Say it once more.
\end{tcolorbox}
\section[\texorpdfstring{\sk{प्रयच्छति} (prayacchati)}{प्रयच्छति (prayacchati)}]{\sk{प्रयच्छति} (prayacchati) — hands over, gives}

\label{lesson:SA-C234-prayacchati}



\begin{quote}
Before the new one: say the Sanskrit for begins, then the Sanskrit for finishes.
\end{quote}

\subsection*{You'll want to know: \sk{प्रयच्छति}}
\textbf{\sk{प्रयच्छति}} — \emph{prayacchati} — \textquotedblleft{}hands over, gives\textquotedblright{}.

\begin{cousinweb}
\emph{prayacchati} is built on \emph{pra}, \textquotedblleft{}forth\textquotedblright{}, and the root √yam, whose present is \emph{yacchati}, \textquotedblleft{}holds\textquotedblright{} (MW). To hold something out towards someone is to hand it over, so MW gives both \textquotedblleft{}stretches forth\textquotedblright{} and \textquotedblleft{}offers, gives\textquotedblright{}.
\end{cousinweb}

\begin{grammarlens}[title={What you've built}]
One more word for this chapter.
\end{grammarlens}

\subsection*{Your turn}
\begin{itemize}
\raggedright
  \item \emph{Say it:} \emph{prayacchati}
  \item \emph{Say it:} \emph{prayacchati}, once more
  \item \emph{Say it:} say \emph{prayacchati}
  \item \emph{Recall:} say the Sanskrit for begins, then the Sanskrit for finishes, then say \emph{prayacchati} again
\end{itemize}

\begin{tcolorbox}[breakable,colback=teal!4,colframe=teal!35!black,title={Before you move on}]
What does \sk{प्रयच्छति} mean? (\textquotedblleft{}Hands over, gives\textquotedblright{}.) Say it once more.
\end{tcolorbox}
